% --- Archon physics formalization source begin ---
% archon:physics
% archon:covers IPhO2026Problems/problem_ipho_2026_e1_b6.lean
% archon:source-report reports/ipho_2026/problem_ipho_2026_e1_b6.source.json
% archon:problem-id ipho_2026_e1
% archon:part-id E1-B6
% archon:previous-part ipho_2026_e1_b5
% archon:previous-part-policy derive_inline_from_problem_only_material

\chapter{Ipho Physics problem ipho\_2026\_e1\_b6}
\label{ch:IPhO2026Problems_problem_ipho_2026_e1_b6}

\paragraph{Problem source.}
\#\# Physical scenario

The inner cylinder contains dry air plus water vapor at total pressure
approximately P\_atm.  The water level is adjusted and its height H is recorded
as temperature T falls.  At T\_0 = 273.15 K, extrapolated height is H\_0 and the
water vapor pressure may be taken as zero.  Vapor pressure obeys
ln(P\_v/P\_v0) = -(Q\_v/R)*(1/T - 1/T\_0).  Use the experimental procedure and
geometry on pages 11--12.

\#\# Current subquestion E1-B6

Convert Q\_v into latent heat per unit mass L\_v and state the formula.

\paragraph{Shared problem context.}
The inner cylinder contains dry air plus water vapor at total pressure
approximately P\_atm.  The water level is adjusted and its height H is recorded
as temperature T falls.  At T\_0 = 273.15 K, extrapolated height is H\_0 and the
water vapor pressure may be taken as zero.  Vapor pressure obeys
ln(P\_v/P\_v0) = -(Q\_v/R)*(1/T - 1/T\_0).  Use the experimental procedure and
geometry on pages 11--12.

\paragraph{Current subquestion.}
Convert Q\_v into latent heat per unit mass L\_v and state the formula.

\paragraph{Figure/image path.}
ipho\_2026\_source/image/E1\_page-12.png

\paragraph{Reusable previous-part conclusions.}
\begin{itemize}
\item Source E1-B5. Question: Construct a Clausius-Clapeyron graph and use it to determine the molar latent heat Q\_v. Policy: derive\_inline\_from\_problem\_only\_material
\end{itemize}

\paragraph{Formalization target.}
create a compiling Lean file with sorry bodies at `IPhO2026Problems/problem\_ipho\_2026\_e1\_b6.lean`.
The Lean declarations must preserve the physical quantities, dimensions or dimensional roles, figure labels, governing-law hypotheses, and final relation expressed by this problem.
Use Mathlib/Physlib names found through LeanExplore where available. If a domain API is missing, introduce faithful local abstractions rather than scalar placeholder aliases.
This is an answer-blind solve-phase task. Derive a candidate result only from the problem-side evidence above; do not assume a target value, select a tolerance around a desired value, or encode a desired result into a candidate domain.
For numerical questions, define the raw end-to-end quantity and apply an explicit source-derived rounding rule. For identification questions, characterize or prove uniqueness before recording the derived candidate.

\begin{theorem}[Ipho Physics formalization target]
\label{thm:physics:ipho_2026_e1_b6:target}
\leanok
The assigned autoformalize agent should translate this IPhO physics problem into Lean declarations in the covered file, with theorem and lemma proof bodies written as `by sorry`.
\end{theorem}
\begin{proof}
\leanok
Fully proved in \texttt{IPhO2026Problems/problem\_ipho\_2026\_e1\_b6.lean}
(namespace \texttt{IPhO2026.E1B6}): the conversion formula
$L_v = Q_v / M_w$ (\texttt{latentHeatPerMass\_eq\_candidate}) follows from the
three defining latent-heat laws on one vaporization sample
($Q = n \cdot Q_v$, $Q = m \cdot L_v$, $m = n \cdot M_w$), which give
$n \cdot Q_v = n \cdot (M_w \cdot L_v)$; cancelling the positive vaporized
amount $n$ yields the bridge form $Q_v = M_w \cdot L_v$
(\texttt{molarLatentHeat\_eq\_molarMass\_mul\_specific}), and dividing by the
positive molar mass $M_w$ yields $L_v = Q_v / M_w$. Positivity of $L_v$ is
certified by \texttt{latentHeatPerMass\_pos}.
\end{proof}
% --- Archon physics formalization source end ---
